\section*{The Propers: Detailed Commentary}

\subsection*{Mercy before the account: the Judge enters judgment\enspace\properrefs{Int., Coll., Gosp.}}

Luke's literal roles remain morally asymmetric: a rich master calls an unjust steward to account. Chrysologus does not erase that level. He moves figural interpretation to a different plane: the rich man is Christ, while the wasteful steward is humanity, like the prodigal son. The account is demanded before death so that there may still be time for confession, satisfaction, and mercy before the final judgment (\emph{Sermon} 125.3--7).

\begin{fourstable}{Role}{Chrysologus's identification}{Action}{Doctrinal force}
\textbf{Rich man and master} & Christ, human in the Incarnation and rich in divinity & Assumes the steward and the steward's prior debts & The creditor is also the rescuing Lord.\\
\textbf{Judge} & Christ & Enters judgment, undergoes penalties, accepts death, and penetrates the underworld & The innocent Judge acts for the accused at his own cost.\\
\textbf{Steward} & Humanity in Sermon 125; the converted Gentile people in Sermon 126 & Squanders, is summoned, returns to itself, and seeks repentance & Redeemed humanity remains the accountable administrator.\\
\textbf{Reduced bonds} & The steward's post-redemption action & Urges repentance, Jubilee mercy, and payment through grace & The rewriting answers Christ's prior payment; it does not constitute it.\\
\end{fourstable}

The central sentence is startlingly compact: \latin{Priora debita suscepi ego, cum te suscepi; solui ego, cum te absolui}---Christ took up the former debts when he took humanity up and paid when he absolved it. Chrysologus then intensifies the exchange: though Judge, Christ stands in judgment; though free from penalty, he bears penalties; the destroyer of death accepts death and the destroyer of hell enters the underworld, not merely to remove punishment but to lead humanity to his dignity (125.8).

\begin{sourcecard}{The critical Latin and Irenaeus}
The critical Latin reads \latin{rei mei reus factus sum}, immediately followed by \latin{a poena liber suscepi poenas}. It describes the Judge taking the accused person's place in judgment: a vicarious pattern that does not make Christ morally sinful or personally liable through guilt. Irenaeus complements it: the incarnate Word, the very one against whom humanity sinned, mediates reconciliation, forgives the debt, and destroys its handwriting on the Cross (\emph{Against Heresies} V.17.1--3). Father and Son are not opposing creditor and deceiver.
\end{sourcecard}

\sectionguard
\subsection*{Debtors made heirs: grace changes the obligation\enspace\properrefs{Ep., Grad., All., Gosp.}}

The strongest verbal bridge comes before the more familiar contrast of children. Romans begins, \latin{debitores sumus, non carni}---we are debtors, but not to the flesh; Luke later names the estate's \latin{debitores domini sui}. Paul is not calling sin a negotiable loan, and Luke is not defining adoption as a contract. Heard together, however, the passages distinguish liberation from irresponsibility: grace breaks the flesh's false mastery so that persons may live under the Spirit as answerable sons and daughters.

\begin{threestable}{Appointed movement}{Direct reception}{Distinction}
\textbf{Debtors, but not to flesh} & Chrysostom treats \latin{moriemini} as a real warning and mortification as death to evil deeds. & The body is God's good work; its substance and ordinary care are not the debt Paul denies.\\
\textbf{Led, yet acting} & Augustine says the Spirit acts in the person who also truly acts; grace is needed for the struggle itself. & Divine agency neither abolishes the will nor excuses passivity.\\
\textbf{Children and heirs} & The Spirit gives filial prayer and witness; Thomas identifies God himself as the inheritance. & Christ is Son by nature; Christians are adopted by grace; Romans 9:4 also names Israel's adoption.\\
\textbf{Children of light} & Luke's rebuke asks whether those who know the final good act with purposeful foresight. & Filial urgency must be purified of the steward's injustice.\\
\end{threestable}

The chants give this freedom its ecclesial posture. The Gradual directs confidence from self to God as refuge; the Alleluia returns the heirs to God's city and holy mountain. Augustine's Psalm 70 exposition makes this transfer from self-confidence to Christ explicit. Thus heirs remain petitioners: costly release creates dependent freedom, not autonomous possession.

\begin{sourcecard}{A hard liturgical boundary --- Romans 8:17a and 17b}
The missal stops after \latin{coheredes autem Christi}. The biblical verse continues \latin{si tamen compatimur, ut et conglorificemur}. Augustine's \emph{Sermon 156} ends at the appointed boundary; Thomas treats the continuation separately. The heir's conformity is therefore cruciform, but suffering neither earns adoption nor becomes a good sought for itself.
\end{sourcecard}

\sectionguard
\subsection*{The rewritten bonds: genuinely different inherited readings\enspace\properrefs{Gosp.}}

Debt imagery by itself cannot determine the steward's identity. The inherited witnesses differ sharply, and their disagreement itself governs how far the image can be pressed.

\begin{fourstable}{Witness}{Steward}{Debts or reductions}{Distinction}
\textbf{Chrysologus 125--126} & Humanity; then the converted Gentile people & Christ pays the former debts; rewriting urges repentance, Jubilee mercy, and grace & Christ is rich master and prior payer, not steward.\\
\textbf{Gaudentius 18} & The devil & Sinners are debtors; reductions are a counterfeit promise that sin will not be charged & Debt allegory here identifies Christ's adversary, not Christ.\\
\textbf{Cyril 108} & No exhaustive assignment & Entrusted wealth is to serve the poor before death and account & Forcing every person, quantity, and commodity obscures the parable.\\
\textbf{Augustine, \emph{Questions} II.34} & The unjust human actor in a contrary comparison & Christian almsgiving should exceed the legal tithe & Not every act is imitable; one must not defraud the Lord in order to give.\\
\textbf{Anthony of Padua} & Every human person & Christ's Passion reconciles humanity before its stewardship is judged & This later saint also makes Christ the rich man, never the steward.\\
\end{fourstable}

Chrysologus's Sermon 126 reclassifies the reductions within its figural reading. The steward, once Gentile and now Christian, returns to the true Lord; the fifty signifies mercy and Jubilee, while the eighty joins the law's ten to grace's octave. His conclusion is explicit: grace pays what nature could not, and the steward sought salvation rather than fraud (126.6--9).

Gaudentius shows what turns on the roles. He calls the devil the unjust steward and treats the reduced debts as a lying assurance that sinners will escape liability. Cyril moves in the other direction and refuses the whole-detail decoding project. Augustine's contrary comparison preserves the literal scandal while extracting moral urgency. No Father or saint among those checked identifies Christ as the unjust steward; Chrysologus and Anthony instead give strong support to the adjacent but more precise claim that Christ is the rich Lord whose costly reconciliation precedes humanity's rewritten account.

\begin{studybox}{Rich Lord, human steward, costly remission}
The patristic witnesses support a rich Lord--human steward--sinful debt--costly remission complex. They do not support transferring the steward's injustice to Christ. Chrysologus's two moments are also distinct: Christ's assumption and payment are prior; repentance, mercy, and grace are figured by the later rewriting.
\end{studybox}

\sectionguard
\subsection*{The unbribable Judge and truthful return\enspace\properrefs{Int., Ep., Gosp., Off., Sec.}}

Augustine gives the indispensable moral reading. The parable is a contrary or a-fortiori comparison, so its every circumstance is not proposed for imitation (\emph{Questions on the Gospels} II.34). Where wealth was acquired by theft or usury, sin must cease and restitution must precede beneficence, as with Zacchaeus; almsgiving cannot bribe Christ the Judge (\emph{Sermon} 113.1--4). Chrysologus produces the deeper paradox: this Judge cannot be bought with stolen gifts, yet freely enters judgment for the accused.

The Offertory sings that God saves the humble and brings proud eyes low. It therefore moves the offering away from the steward's self-protective control. The Secret completes the correction: worshippers truly bring gifts, but the gifts come \latin{de tua largitate}; grace's operative power sanctifies conduct and leads it toward everlasting joy. Irenaeus likewise says that the Church offers what is already God's, while an outward sacrifice cannot conceal injustice within (\emph{Against Heresies} IV.18.5--6).

\begin{fourstable}{Verbal bond}{One use}{Other use}{Theological work}
\latin{debitores} & The master's debtors owe within the parable & Christians owe nothing to the flesh's mastery & Redemption changes servitude without abolishing answerability.\\
\latin{redde / reddit} & The steward must render an account & The Spirit renders witness to filial identity & Judgment exposes administration; grace attests adoption.\\
\latin{Suscepimus / Suscipe} & The Church has received mercy & The Church asks God to receive his own gifts & Prior gift becomes thankful return, never barter.\\
\latin{cautio / chirographum} & Luke's written bond is reduced & Patristic reception evokes sin's canceled handwriting & A documentary image serves repentance and the Cross without making salvation a falsified ledger.\\
\end{fourstable}

Luke's canonical continuation keeps the poor and final judgment in view. Mammon can become a rival master; the rich man ignores Lazarus at his gate. Yet the needy are not independent brokers selling heaven. Mercy toward Christ in them is the graced fruit of reconciliation and truthful stewardship, not its purchase price.

\sectionguard
\subsection*{From canceled handwriting to the restored body\enspace\properrefs{Ep., Gosp., Comm., Postcomm.}}

Origen's \emph{Homily on Genesis} 13.4 directly invokes Luke's rewritten bills as an image of a wider mystery. The documents of sin are ours; at the Cross the handwriting is fastened to Christ's Cross, and in baptism it is blotted out in the font. Origen does not identify the steward with Christ. He takes the documentary image beyond the steward's act to the Paschal and sacramental destruction of sin's bond.

Irenaeus supplies the canonical-doctrinal depth. The incarnate Mediator cancels disobedience by obedience, forgives the debt because humanity sinned against him, and destroys the handwriting by fastening it to the Cross (\emph{Against Heresies} V.17.1--3). Later he argues that Christ's Eucharistic cup and bread nourish flesh that, though buried and dissolved, will be raised by the Word to incorruption (V.2.2--3). Cancellation is therefore ordered not merely to a cleared account but to communion and resurrection.

This clarifies the Epistle's anthropology. Paul commands death to the \emph{deeds} of the flesh, not hatred of bodily substance. The Communion then commands \latin{Gustate et videte}; Augustine reads the verse of Christ's Flesh and Blood. Finally the Postcommunion asks \latin{reparatio mentis et corporis}. The body whose disordered deeds are mortified is the body Christ nourishes and promises to restore.

\begin{threestable}{Sacramental claim}{What follows}{What does not follow}
\textbf{Christ acts} & The sacramental effect is objective before it is emotionally perceived (CCC 1127--1129). & Reception is not mechanical or independent of conversion.\\
\textbf{The communicant receives} & Worthy reception nourishes charity and orders embodied life toward glory. & Felt sweetness is not the sole measure of fruit, and grave sin is not bypassed.\\
\textbf{Mind and body are restored} & Present grace reaches conduct; resurrection is bodily consummation. & The prayer promises neither instant cure nor glory detached from conformity to Christ.\\
\end{threestable}

The appointed chants now disclose one response rather than a separate theme: mercy is received within God's just city; heirs seek refuge rather than self-possession; proud eyes are humbled at the offering; and Eucharistic tasting returns the redeemed steward to hope in the Lord. The sequence ends neither with a manipulated account nor with autonomous solvency, but with forgiven persons receiving the Lord who first received them.

\sectionguard
\section*{The Propers: Notable and Quotable}

{\small
\begin{itemize}[itemsep=0.35em,topsep=0.2em,parsep=0pt]
\item \textbf{Gospel, Luke~16:2--4 --- Bill Nye, ``The Parable of the Unjust Steward,'' in \emph{Bill Nye and Boomerang} (1883).} ``Give an account of thy stewardship'' is addressed to Keno-El-Pharo, a corrupt hotel clerk whose debtors rewrite bar and lodging bills in dollars, replacing Luke's oil and wheat accounts. He cannot obtain a ``position as mining expert'' and finally lands in Kansas City. The mock-biblical Western relocates the entire parable into hotel, gambling, and mining culture while preserving its self-saving fraud.

\item \textbf{Gospel, Luke~16:3 --- Oliver Wendell Holmes, \emph{The Poet at the Breakfast-Table}, sec.~II (1872).} Holmes applies ``I cannot dig, to beg I am ashamed'' to a genteel woman trained to be ornamental rather than useful and stranded when fortune recedes. The steward's panicked excuse becomes a social diagnosis: economic helplessness can be manufactured by class education, so the irony falls partly on the culture that prepared her for display.

\item \textbf{Gospel, Luke~16:9 --- E.~F. Benson, \emph{Mammon and Co.}, book~I, ch.~I, ``The City Dinner'' (1899).} Benson calls fashionable insiders ``the mammon of unrighteousness'' and coins their comic opposite, ``the unmammon of righteousness.'' Both groups sustain Lord Conybeare's profitable duplicity---one knowingly, the other through pious credulity---so suspect wealth becomes a social caste and the invented antonym sharpens the satire.
\end{itemize}
}

\sectionguard
\section*{The Propers: Interpretive Possibilities}

\noindent\emph{The following exploratory proposals were developed in an AI-assisted editorial process. They are attributed to none of the cited authorities and claim no historical compositional intent. Each states its anchors and limits.}

\Needspace{8\baselineskip}
\noindent\textbf{Mercy arrives before the books open.\enspace\properrefs{Int., Coll., Gosp.}}
The Introit receives mercy in the Temple and names God's justice; the Collect asks right thought and action; only then does the Gospel demand an account. This order joins prior mercy, enabled agency, and accountability. Alone, mercy can become indulgence, dependence passivity, and judgment panic; together, the audit occurs inside mercy's power for change. Its fruit is urgency without presumption or despair. Liturgical sequence is not Luke's chronology, and prior mercy removes neither restitution, conversion, nor judgment.

\Needspace{8\baselineskip}
\noindent\textbf{Christ is the steward's costly inverse.\enspace\properrefs{Gosp., Sec., Comm.}}
The steward reduces claims to purchase reception; Chrysologus makes Christ rich Lord and prior payer. The Secret returns gifts from bounty; Communion gives the Lord, not a manipulated asset. Reversed ownership, cost, and hospitality answer self-preservation with truthful return and self-gift---a relation the Gospel alone cannot supply. Its fruit is transformed stewardship without making fraud an image of the Cross. No checked Father calls Christ the unjust steward; the inverse is editorial, not the parable's key.

\Needspace{8\baselineskip}
\noindent\textbf{The canceled bond grows a restored body.\enspace\properrefs{Gosp., Comm., Postcomm.}}
Luke supplies rewritten bills; Origen and Irenaeus redirect handwritten debt to the Cross; Communion moves from paper to tasting; the Postcommunion asks embodied restoration. The mechanism passes from inscription through reception to repair. Alone, forgiveness can remain external and Eucharistic restoration unrelated; together they order acquittal to action and resurrection. Its fruit refuses opposition between forensic and participatory salvation. Luke's steward falsifies accounts whereas Christ truthfully destroys sin's bond; the analogy guarantees no immediate felt or moral change.

\Needspace{8\baselineskip}
\noindent\textbf{An heir can remain a steward.\enspace\properrefs{Ep., Gosp., Sec.}}
Romans denies the flesh's claim while naming heirs; Luke's steward must account; the Secret returns gifts from God's bounty. Transformed possession ends servile mastery without creating autonomous ownership. Alone, heirship becomes entitlement and stewardship insecurity; together, belonging sharpens answerability. Its fruit is a grammar for authority and wealth: receive without servility, act without possessiveness, account without manipulation. Paul's debt, Luke's liabilities, and the Secret's bounty differ; adoption is neither earned by management nor humanly revocable.

\Needspace{6\baselineskip}
\noindent\textbf{Eternal dwellings turn assets into hospitality.\enspace\properrefs{Int., All., Gosp., Sec.}}
The Introit receives mercy in the Temple and praises God's city; the Alleluia returns there before Luke's eternal dwellings; the Secret points toward everlasting joy. Space enlarges from sanctuary through city to eschatological reception, turning entrusted goods into communion. Alone, city praise, almsgiving, and joy remain inert; together, salvation's ``where'' is receiving and being received. Its fruit is hospitality without leverage. Alms cannot purchase heaven, ``friends'' do not replace God, and the liturgical city is not proved Luke's referent.
